% Shared implementation section for the general and WP7 Feel++ fiches.
% Render once if both fiches are included in the same document.
\ifdefined\FeelppImplementationTwentySixRendered\else
\global\let\FeelppImplementationTwentySixRendered\relax
\newcommand{\FeelppImplementationAdvance}[1]{%
  \ifcsname r@tab:advance:#1\endcsname
    \Cref{tab:advance:#1}%
  \else
    \nolinkurl{#1}%
  \fi}

\subsection{2026 implementation developments}
\label{sec:Feelpp:implementation-2026}

The following developments implement the Cemosis contributions to WP1,
WP3, WP4, WP5 and WP6 in the Feel++ framework. Scientific WP chapters
present the corresponding methods; this section identifies their software
components, examples and verification sources.

\subsubsection{Solid-shell elements and tensor operators}
\label{sec:Feelpp:implementation-2026:sb9}
\href{https://github.com/feelpp/feelpp/pull/2510}{Feel++ PR~2510} supplies SB9 shell geometry and strain
operators, symmetric tensor representations in Voigt and Mandel notation,
constitutive contractions, assumed-shear and stabilization operators.
Shell geometry is cached for reuse by the membrane--bending, pinching and
shear terms. The mixed formulation uses a continuous vector $Q_1$
displacement space and a discontinuous $P_0$ internal cell mode.
The implementation maintains explicit component and tensor-storage
orderings and supports block-aware assembly.

Verification configurations are supplied under
\nolinkurl{feelpp/quickstart/elasticity/shell}, including bending-patch,
plate and SB9 validation cases. Expression and tensor microbenchmarks are
provided under \nolinkurl{benchmarks/feelpp}. The discretization contribution
is described in \FeelppImplementationAdvance{WP1-A2026-2}.

\subsubsection{Distributed geometric search}
\label{sec:Feelpp:implementation-2026:bvh}
\href{https://github.com/feelpp/feelpp/pull/2494}{Feel++ PR~2494} refactors parallel BVH intersection queries
using a two-level hierarchy. Communications use nonblocking MPI operations,
bounded batches and an event loop with a minimum interval between
communication pulses. The two-dimensional implementation shares the
three-dimensional query machinery. The documented \texttt{anyHint} and
\texttt{closest} contexts are operational; the all-intersections context
can return incomplete results in parallel. This scope is retained in
\FeelppImplementationAdvance{WP1-A2026-3}.

\subsubsection{HPDDM integration and block static condensation}
\label{sec:Feelpp:implementation-2026:solvers}
The HPDDM interface in \nolinkurl{feelpp/feel/feelalg/preconditionerpetsc.hpp}
adds HPDDM preconditioner configuration through PETSc, including auxiliary
Neumann matrices and index sets. The solver interface is exercised by
\nolinkurl{testsuite/feelalg/test_hpddm.cpp},
\nolinkurl{testsuite/feelalg/test_hpddm_laplacian.cpp} and elasticity
quickstart tests.

The component \nolinkurl{feelpp/feel/feelalg/staticcondensation.hpp}
extends static condensation to two-by-two block systems, with product-space
tests. Mixed assembly and row-wise Dirichlet elimination are supported by
\href{https://github.com/feelpp/feelpp/pull/2437}{Feel++ PR~2437}.
These components realize the solver contribution described in
\FeelppImplementationAdvance{WP3-A2026-2}; they do not replace the separately documented
HPDDM framework contribution or its existing benchmark results.

\subsubsection{Variational optimization through PETSc TAO}
\label{sec:Feelpp:implementation-2026:tao}
\href{https://github.com/feelpp/feelpp/pull/2504}{Feel++ PR~2504} introduces a Feel++ interface to PETSc TAO.
The API connects objectives, gradients, Hessians, constraints and bounds
to solver configuration, monitoring and convergence reporting.
Functional and variational-form tests exercise the interface.
The implementation supports the deterministic inverse-problem role in
\FeelppImplementationAdvance{WP4-A2026-2} and the continuous-optimization role in
\FeelppImplementationAdvance{WP5-A2026-4}, using one shared software interface.

\subsubsection{preCICE coupling and OpenFOAM deployment}
\label{sec:Feelpp:implementation-2026:precice}
The preCICE integration identified by
\href{https://github.com/feelpp/feelpp/issues/2525}{Feel++ issue~2525}
provides coupling mesh and field
interfaces. The adapter supports multiple coupling meshes and field
sampling at element barycentres, including discontinuous spaces.
The diffusion quickstart under \nolinkurl{feelpp/quickstart/precice}
uses nonmatching trace meshes and exchanges temperature between two
participants. The receiving solver applies the transferred field through
Feel++'s variational Dirichlet interface. The example's analytical solution,
acceptance criterion and MPI configurations are summarized in
\FeelppImplementationAdvance{WP3-A2026-3}. It performs one-way temperature transfer;
heat-flux feedback and iterative Dirichlet--Neumann coupling are outside
this example.

The recipes under \nolinkurl{packaging/docker/openfoam} supply a separate
OpenFOAM participant image with preCICE and its adapter, together with
Docker and Apptainer/Slurm launch examples. Each participant retains its
own MPI installation; preCICE sockets exchange coupling data. Dependency
versions, adapter archives and image artifacts are identified by checksums.
Runtime smoke tests check dependency versions, dynamic-library resolution,
adapter loading and MPI execution. These deployment checks are distinct
from numerical validation of a coupled OpenFOAM--Feel++ application.

\subsubsection{Streaming statistics and quantile sketches}
\label{sec:Feelpp:implementation-2026:stats}
The module \nolinkurl{feelpp/feel/feelstats}
adds \texttt{feelstats}: online mean and variance, extrema,
paired covariance/correlation, central moments through fourth order,
and a mergeable KLL-style quantile sketch. Statistical accumulators
and sketches support serialization and merging of partial summaries.
Tests under \nolinkurl{testsuite/feelstats} cover statistical summaries,
merging, quantile selection and mixed-scale samples.
The sketch uses equal per-level capacities and explicitly does not claim
Apache DataSketches' normalized-rank error bounds.
These components implement the WP6 statistical-processing contribution in
\FeelppImplementationAdvance{WP6-A2026-2} and support timing diagnostics.

Benchmark results are reported in the corresponding application fiches in \Cref{chap:applications}.
\fi
